\section{Discussion}
\label{sec:discussion}

\paragraph{What does healing restore?}
The combined evidence argues against treating recovery as a single scalar. PPL
measures how well the healed model fits held-out next-token events; it improves
substantially in every inherited-prefix arm. That improvement is useful---the
models recover fluent continuation and several commonsense tasks---but it does
not uniquely determine answer selection or closed-book recall. The endpoint
scatter makes the distinction concrete: models with similar PPL can occupy
meaningfully different MMLU regimes, and random initialization can outrank a
frozen inherited front in PPL while underperforming it on answer letters. We
therefore interpret healing as partial recovery of distributional modeling plus
a task-dependent remainder, not as a binary return of the original model.

\paragraph{Interface sensitivity is substantial but insufficient.}
Complete-option MMLU recovers more smoothly than answer-letter MMLU, showing
substantial protocol sensitivity that is consistent with an answer-symbol
component. The protocols also change the prompt, candidate continuation,
tokenization, and normalization, so they do not isolate mapping loss. The
random-init fluency floor and the independent PopQA, TriviaQA,
and NQ-open gaps rule out a formatting-only account. This distinction matters
for evaluation: a single multiple-choice protocol can overstate damage when its
interface is brittle, whereas a content-likelihood protocol can overstate
recovery when generic fluency makes plausible options easy to score. Agreement
across interfaces is more informative than designating either as ground truth.

\paragraph{Different constructions yield different endpoints.}
The step-matched keep14--ShortGPT comparison shows that half depth is not one
operating point. Preserving selected upper blocks yields a much stronger
likelihood--capability endpoint than retaining only a prefix and regrowing a
tail. The result does not identify whether the advantage comes from the
original final block, two additional inherited blocks, non-contiguous selection,
or avoiding fresh layers. It shows that conclusions drawn from one prefix
construction should not be generalized to every 16-layer construction;
isolating the responsible factor requires additional matched structures.

\paragraph{Which behaviors remain difficult?}
The residual deficit spans factual recall (PopQA, TriviaQA, and NQ-open), broad
academic knowledge (all four MMLU groups), and some answer-interface behavior.
In contrast, several continuation-style commonsense tasks recover more readily.
This taxonomy is behavioral rather than anatomical. Knowledge-neuron and causal
tracing studies motivate asking whether particular feed-forward computations
support factual recall~\citep{dai2022knowledge,meng2022locating}, but neither our
subject groups nor ShortGPT's retention of block 31 identifies the missing
substrate. Structure-isolation and causal interventions are required before
making that claim.

\paragraph{A reporting protocol for structural compression.}
For prune-then-heal studies, we recommend reporting likelihood, target
capability, recovery compute, and exact structural policy together. An intact
model continued on the same corpus and nominal schedule measures short-horizon
drift, while same-shape
initialization/adaptation controls expose fluency floors and inherited-computation
effects.

\begin{center}
\fbox{\begin{minipage}{0.92\columnwidth}
\small\textbf{Before claiming recovery from PPL alone:}
(1) evaluate an independent target capability;
(2) test a second scoring or generation interface;
(3) include a continued-full-model corpus control; and
(4) report the exact structure and healing budget.
\end{minipage}}
\end{center}

This checklist does not replace efficiency or endpoint-quality comparisons. It
prevents a low PPL alone from being interpreted as evidence that the behaviors
of interest have returned.
